\documentclass[10pt,a4paper]{amsart}
\usepackage[margin=19mm]{geometry}
\usepackage{amsmath,amssymb,mathtools}
\usepackage[scr=rsfs,cal=euler]{mathalfa}
\usepackage{xfrac}
\usepackage[unicode,hidelinks,bookmarks=false]{hyperref}
\newcommand{\ie}{i.e.\ }
\newcommand{\cf}{cf.\ }
\newcommand{\resp}{respectively\ }
\newcommand{\Subsection}{\subsection}
\providecommand{\nobreakdash}{\nobreak\hbox{-}}
\setlength{\emergencystretch}{2em}
\multlinegap0pt
\usepackage{bm}
\usepackage{bbm}
\usepackage{dsfont}
\usepackage{xspace}
\usepackage{cancel}
\usepackage{mathtools}
\usepackage{amsthm}
\makeatletter
\@ifundefined{lemma}{\newtheorem{lemma}{Lemma}}{}
\@ifundefined{proposition}{\newtheorem{proposition}{Proposition}}{}
\makeatother
\def\Q{\mathbb Q}
\def\Z{\mathbb Z}
\newcommand{\F}{\mathbb{F}}
\newcommand{\Gal}{\mathrm{Gal}}
\newcommand{\dlog}{\mathrm{dlog}}
\newcommand{\nr}{\mathrm{nr}}
\newcommand{\ord}{\mathrm{ord}}
\newcommand{\val}{\mathrm{val}}
\title[$p$-adic approaches to unlikely intersections]{$p$-adic approaches to unlikely intersections}
\author{Netan Dogra}
\date{Source: arXiv:2504.10611v2 (CC BY 4.0)}
\begin{document}
\maketitle

\noindent\emph{Source and licence.} $p$-adic approaches to unlikely intersections. Version arXiv:2504.10611v2 (CC BY 4.0), distributed under the Creative Commons Attribution 4.0 International licence, as recorded in the arXiv licence metadata for this exact version and shown on the abstract page https://arxiv.org/abs/2504.10611v2. Full rights notice: licenses/arxiv-2504.10611-CC-BY-4.0.txt.\par\smallskip

\noindent\textit{Editorial scope.} Counted unit: the Proposition whose source label is \texttt{prop:ram} in the article (source0.tex lines 236--238, source label \texttt{prop:ram}), delivered together with the article's own complete original proof of it (source0.tex lines 239--260, one \texttt{proof} environment of 22 lines). No mathematics is written, repaired or re-derived: statement and proof are the article's own text line for line. Required context retained uncounted: lemma:slop (lines 186--193); lemma:integral (lines 202--211). Every definition, display and label of those lines is retained unchanged. Omitted from this package and not redistributed: 1--185, 194--201, 212--235, 261--435. Nothing inside a retained excerpt is omitted or altered: every retained body is delivered exactly as the article's lines are written, so no omission receipt of any kind accompanies this package. Citations: the retained excerpts carry no citation command at all, so no bibliography entry of the article is transcribed and no reference is redistributed. Reference adaptation: none was needed and none was made; every reference inside the delivered unit resolves inside it, so no unresolved reference or citation is produced. Printed slips kept literal: no printed word, symbol, hypothesis or number of the counted statement or of its proof was altered. Layout and class: the article's own class is \texttt{amsart} with packages including amsmath, amsfonts, amssymb, amsthm, enumerate, tikz-cd, color; the wrapper is the lane's pinned \texttt{amsart} editorial wrapper at 10pt on A4, which restates the theorem-like environments the retained text needs and adds the article's own macro definitions that the retained text needs (\texttt{\textbackslash F}, \texttt{\textbackslash Gal}, \texttt{\textbackslash Q}, \texttt{\textbackslash Z}, \texttt{\textbackslash dlog}, \texttt{\textbackslash nr}, \texttt{\textbackslash ord}, \texttt{\textbackslash val}), with extra wrapper packages (bm, bbm, dsfont, xspace, cancel, mathtools). The delivered numbering is the unit's own. Assets: no retained span contains a figure environment, a graphics-inclusion command, a PDF-page inclusion, a table or a TikZ picture; no image file is copied into this package, the package declares no asset dependency and no image file is redistributed with it. Licence: version arXiv:2504.10611v2 of this article is distributed under the Creative Commons Attribution 4.0 International licence, as recorded in the arXiv licence metadata for this exact version and shown on the abstract page https://arxiv.org/abs/2504.10611v2. A full rights notice is delivered at licenses/arxiv-2504.10611-CC-BY-4.0.txt. Characters: every retained excerpt is pure ASCII, so the delivered engine, XeLaTeX with its default Latin Modern text font, reports no missing character.
\begin{lemma}\label{lemma:slop}
Let $k =\ord _{\overline{z}}(\frac{df}{f})+1.$ Suppose $k<p$.
\begin{enumerate}
\item If $\val (\log (f(z_0 )))>0$, then the negative slopes $\lambda$ of the Newton polygon of $\log (f)$ to the right of $k$ are of the form $\frac{1}{k(p^n-p^{n-1})}$, for $n>0$. The corresponding end vertices are $(kp^n,-n)$ for $n\geq 0$.
\item If $\val (\log (f(z_0 )))=v\leq 0$, the negative slopes $\lambda $ of the Newton polygon of $\log (f)$ are $\frac{1}{k p^{v+1}}$ and $\frac{1}{k(p^{n+1}-p^{n})}$ for $n>v$. The corresponding end vertices are $(0,0)$ and $(kp^n ,-n)$ for $n>v$.
\item If $\log (f)$ contains a non-integral negative slope $>\frac{1}{p-1}$, then $\frac{df}{f}$ vanishes at $\overline{z}$ (i.e. $k>1$).
\end{enumerate}
\end{lemma}

\begin{lemma}\label{lemma:integral}
\begin{enumerate}
\item Let $X$ be a geometrically irreducible curve over $\Z _p ^{\nr}$ with a map $(f_1 ,\ldots ,f_n ):X\to \mathbb{G}_m ^n $. Let $C$ be the projective closure of the normalisation $\widetilde{X}$ of $X$ and $\{ P_1 ,\ldots ,P_k \}=C-\widetilde{X}$ be the union of the divisors of the $f_i$. Suppose the $P_i$ are pairwise disjoint modulo $p$ (i.e. the divisor $\sum P_i $ is \'etale over $\Z _p ^{\nr}$). Then there are are finitely many surjective homomorphisms
\[
\theta _j :\mathbb{G}_m ^n \to \mathbb{G}_m ^{n-1}
\]
such that if $z\in X(\Q _p ^{\mathrm{nr}})\cap (\mathbb{G}_m ^n )^{[2]}$, and $f_i (z)$ is not in $R^\times $, then the image of $z$ in $(\mathbb{G}_m )^{n-1}$ under one of the $\theta _j$ lies in $(\mathbb{G}_m ^{n-1})^{[2]}(R)$. In particular, when $n=3$ all but finitely many points of $X(\overline{\Q }_p )\cap (\mathbb{G}_m ^3 )^{[2]}$ are integral.
\item Let $X$ be a geometrically irreducible curve over the $S$-integers of a number field $K$, and let $(f_1 ,\ldots ,f_n ):X\to \mathbb{G}_m ^3$ be an $\mathcal{O}_{K,S}$-morphism. Then there is a finite set $T\supset S$ of primes of $K$ such that all but finitely many $\overline{K}$-points of $X(\overline{K})\cap (\mathbb{G}_m ^3 )^{[2]}$ are $T$-integral.
\end{enumerate}
\end{lemma}

\begin{proposition}\label{prop:ram}
Let $X$ be a smooth projective curve over $\Z _p ^{\nr }$ of genus $g$ and $f_1 ,\ldots ,f_n$ functions with pairwise disjoint divisors modulo $p$. Let $Z$ be the union of all points where one of the $f_i$ has a zero or pole, and let $d$ be the degree of the divisor $Z$. If $p\geq 2g+d$ and the $\overline{\Q }_p $-points of $Z$ have pairwise distinct reductions in $\overline{\F }_p$, then every point of $X(\overline{\Q }_p )\cap (\mathbb{G}_m ^{n})^{[2]}$ has ramification degree at most $2g+d$.
\end{proposition}

\begin{proof}
By Lemma \ref{lemma:integral}, we may reduce to the case where all points on $X(\overline {\Q }_p )\cap (\mathbb{G}_m ^{r})^{[2]}$ are integral with respect to the $f_i$. Let $z$ be a ramified point in $X(\overline{\Q }_p )\cap (\mathbb{G}_m ^r )^{[2]}$. Let $\overline{z}$ be the reduction of $z$ modulo $p$. We may choose a point $z_0 \in X(\Q _p ^{\mathrm{nr}})$ congruent to $z$ modulo $p$, and an integral parameter $T$ at $z_0$, such that $\val (T(z))$ is not in $\Z $.

Recall that every codimension $d$ connected subgroup of $\mathbb{G}_m ^n $ is the connected component of the kernel of a surjection $\mathbb{G}_m ^n \to \mathbb{G}_m ^d $. Let $g_1 $ and $g_2 $ be composite maps
\[
X\to \mathbb{G}_m ^r \to \mathbb{G}_m ^2
\]
such that $g_i (z)$ is torsion. If $]\overline{z}[$ contains a point in $(\mathbb{G}_m ^n )^{[2]}$ of ramification degree $N$, then it contains at least $N$ such points, since $X(\overline{\Q }_p )\cap (\mathbb{G}_m ^n )^{[2]}$ is stable under the action of $\Gal (\overline{\Q }_p |\Q _p ^{\nr})$. Hence it is enough to prove that the number of common slopes of $g_1 $ and $g_2$, counted with multiplicity, is less than $2g+d$. By the description of the slopes of $\log (g_i )$ in Lemma \ref{lemma:slop}, it is enough to prove that, for any $\lambda <\frac{1}{2g+d}$, there exist $\mu _1 ,\mu _2 \in \Z$ such that $\lambda $ is not a slope of $\mu _1 \log (g_1 )+\mu _2 \log (g_2 )$.

Since all positive slopes less than $\frac{1}{2g+d}$ are of the form $\frac{1}{kp^i}$ or $\frac{1}{k(p^i -p^{i-1})}$, and $k<2g+d<p$, it is enough to first prove the result for $\lambda $ of the form $\frac{1}{k(p^i-p^{i-1})}$ and then for $\lambda $ of the form
Without loss of generality we may assume $\ord _{\overline{z}}\frac{dg_1 }{g_1 }\leq \ord _{\overline{z}}\frac{dg_2 }{g_2 }$. Define $\omega _1 =\dlog (g_1 )$ and $\omega _2 =\dlog (g_2 )-\alpha \dlog (g_1 )$, where $\alpha$ is the value of $\frac{\dlog g_2 }{\dlog (g_1 )}$ at $\overline{z}$. Then 
\[
\int ^z \omega _1 =\int ^z \omega _2 =0.
\]
Let $k_i $ be the order of vanishing of $\omega _i$ at $\overline{z}$, and let $a_i$ be the valuation of $\int ^{z_0 }_b \omega _i $. Then $k_1 <k_2 \leq 2g+d$. On the other hand by Lemma \ref{lemma:slop}, the slopes of $\int \omega _i $, viewed as power series in $T$, are $\frac{1}{k_i p^{a_i +1}}$ and/or $\frac{1}{k_i (p^n-p^{n-1})}$. Since $p>k_2 >k_1 $, $\int \omega _1 $ and $\int \omega _2 $ do not have a common (negative) slope less than $\frac{1}{2g+d}$ of the form $\frac{1}{k_i (p^n-p^{n-1})}$. 

We now consider slopes of the form $\frac{1}{kp^i}$. Let $v_i$ be the valuation of $\log (g_i (z_0 ))$. If $v_1 =v_2$, we may similarly we may choose $\mu _1 ,\mu _2 \in \Z -p\Z $ such that 
\[
\val (\mu _1 \log (g_1 (z_0 ))+\mu _2 \log (g_2 (z_0 )))>v_1 .
\]
It follows that if $\log (g_1 )$ and $\log (g_2 )$ both have a slope of the form $\frac{1}{k_i p^{v_i }}$, then this is not a slope of $\mu _1 \log (g_1 )+\mu _2 \log (g_2 )$.
\end{proof}

\end{document}
